\documentclass[a4paper,12pt]{article}
\usepackage[utf8]{inputenc}
\usepackage[polish]{babel}
\usepackage[T1]{fontenc}
\usepackage{hyperref}
\usepackage{listings}
\usepackage{xcolor}
\usepackage{geometry}

\geometry{margin=2.5cm}

% --- Konfiguracja wyglądu kodu (Python) ---
\definecolor{codegreen}{rgb}{0,0.6,0}
\definecolor{codegray}{rgb}{0.5,0.5,0.5}
\definecolor{codepurple}{rgb}{0.58,0,0.82}
\definecolor{backcolour}{rgb}{0.95,0.95,0.92}

\lstdefinestyle{pythonstyle}{
    backgroundcolor=\color{backcolour},   
    commentstyle=\color{codegreen},
    keywordstyle=\color{magenta},
    numberstyle=\tiny\color{codegray},
    stringstyle=\color{codepurple},
    basicstyle=\ttfamily\footnotesize,
    breakatwhitespace=false,         
    breaklines=true,                 
    captionpos=b,                    
    keepspaces=true,                 
    numbers=left,                    
    numbersep=5pt,                  
    showspaces=false,                
    showstringspaces=false,
    showtabs=false,                  
    tabsize=4
}
\lstset{style=pythonstyle}

\title{Dokumentacja Biblioteki GUI\\ (E-Trener Framework)}
\author{Dokumentacja Techniczna}
\date{\today}

\begin{document}

\maketitle
\tableofcontents
\newpage

\section{Wprowadzenie}
Niniejsza dokumentacja opisuje prostą, reużywalną bibliotekę GUI opartą na architekturze \textbf{PySide6}. Framework został zaprojektowany z myślą o aplikacjach posiadających sztywny, dwukolumnowy układ:
\begin{itemize}
    \item \textbf{Lewa kolumna:} Dynamiczne menu boczny z opcjami (przyciskami).
    \item \textbf{Prawa kolumna:} Główny panel zawartości (np. wykresy, formularze, odtwarzacz wideo).
\end{itemize}

Biblioteka eliminuje potrzebę zamykania i otwierania nowych okien podczas nawigacji. Przełączanie między ekranami odbywa się wewnątrz jednego okna za pomocą mechanizmu \texttt{QStackedWidget}. Pasek podziału został zablokowany, dzięki czemu interfejs zachowuje stałe proporcje szerokości.

\section{Główne Klasy API}

\subsection{Klasa \texttt{Screen}}
Klasa reprezentuje pojedynczy stan (widok) aplikacji. Składa się z kontenera na przyciski menu oraz przypisanego do nich widżetu zawartości.

\subsubsection*{Konstruktor}
\begin{lstlisting}[language=Python]
Screen(content_widget: QWidget)
\end{lstlisting}
\textbf{Argumenty:}
\begin{itemize}
    \item \texttt{content_widget} -- Dowolny obiekt klasy pochodnej od \texttt{QWidget}, który ma zostać wyświetlony po prawej stronie okna (np. panel z wykresem).
\end{itemize}

\subsubsection*{Metody}
\begin{itemize}
    \item \texttt{add_option(text: str, callback: function)} -- Dodaje przycisk do menu bocznego dla tego konkretnego ekranu.
    \begin{itemize}
        \item \texttt{text} -- Napis wyświetlany na przycisku.
        \item \texttt{callback} -- Referencja do funkcji języka Python (np. \texttt{lambda} lub metoda klasy), która zostanie wykonana po kliknięciu.
    \end{itemize}
\end{itemize}

\subsection{Klasa \texttt{AppWindow}}
Główne okno aplikacji, pełniące rolę silnika renderującego i menedżera widoków.

\subsubsection*{Konstruktor}
\begin{lstlisting}[language=Python]
AppWindow(title="Aplikacja", bg_image_path="gui_background.jpeg")
\end{lstlisting}
\textbf{Argumenty:}
\begin{itemize}
    \item \texttt{title} -- Tekst wyświetlany na pasku tytułowym okna.
    \item \texttt{bg_image_path} -- Ścieżka do pliku graficznego tła (automatycznie rozmywanego efektem blur).
\end{itemize}

\subsubsection*{Metody publiczne}
\begin{itemize}
    \item \texttt{register_screen(name: str, screen: Screen)} -- Rejestruje obiekt ekranu pod unikalną nazwą tekstową i automatycznie generuje dla niego strukturę widżetów Qt.
    \item \texttt{switch_to_screen(name: str)} -- Natychmiastowo podmienia zawartość lewego menu oraz prawego panelu na te zdefiniowane w ekranie o nazwie \texttt{name}.
    \item \texttt{close()} -- Zamyka okno aplikacji (wbudowana metoda Qt, idealna jako callback dla przycisku "Wyjście").
\end{itemize}

\newpage
\section{Przykład Użycia (Quick Start)}

Poniższy kod prezentuje minimalny, działający przykład implementacji dwóch ekranów z możliwością płynnej nawigacji między nimi.

\begin{lstlisting}[language=Python, caption=Przykład implementacji aplikacji za pomocą biblioteki]
import sys
from PySide6.QtWidgets import QApplication, QLabel, QWidget, QVBoxLayout
from PySide6.QtCore import Qt
from gui_lib import AppWindow, Screen 

def create_page(text):
    """Funkcja pomocnicza tworząca prosty widżet zawartości."""
    page = QWidget()
    layout = QVBoxLayout(page)
    label = QLabel(text)
    label.setStyleSheet("color: white; font-size: 22px;")
    label.setAlignment(Qt.AlignCenter)
    layout.addWidget(label)
    return page

if __name__ == "__main__":
    app = QApplication(sys.argv)
    
    # 1. Inicjalizacja głównego okna
    window = AppWindow(title="E-Trener v2")

    # 2. Definicja Ekranu 1 (Startowego)
    start_content = create_page("Witaj w aplikacji!\nWybierz opcję z menu.")
    screen_start = Screen(start_content)
    screen_start.add_option("Panel Treningów", lambda: window.switch_to_screen("treningi"))
    screen_start.add_option("Wyjście", window.close)
    window.register_screen("start", screen_start)

    # 3. Definicja Ekranu 2 (Menu Treningów)
    training_content = create_page("Sekcja Treningowa\n[Miejsce na Wykres/Wideo]")
    screen_training = Screen(training_content)
    screen_training.add_option("Dodaj serię", lambda: print("Seria dodana!"))
    screen_training.add_option("<- Powrót do menu", lambda: window.switch_to_screen("start"))
    window.register_screen("treningi", screen_training)

    # 4. Uruchomienie aplikacji
    window.switch_to_screen("start")
    window.show()
    sys.exit(app.exec())
\end{lstlisting}

\section{Skalowanie i Wygląd}
\begin{itemize}
    \item \textbf{Responsywność:} Okno uruchamia się domyślnie w rozmiarze stanowiącym $80\%$ rozdzielczości ekranu użytkownika, zachowując minimalny rozmiar bezpieczny ($900 \times 600$ pikseli). Dynamiczne marginesy są przeliczane podczas zdarzenia \texttt{resizeEvent}.
    \item \textbf{Proporcje:} Podział okna jest ustalony w stosunku \textbf{3:7} (30\% szerokości dla menu bocznego, 70\% dla panelu treści).
    \item \textbf{Style CSS:} Wygląd (m.in. efekt \textit{hover} na przyciskach, kolory rgba oraz zaokrąglenia krawędzi) zdefiniowany jest globalnie w arkuszu stylów wewnątrz klasy \texttt{AppWindow}.
\end{itemize}

\end{document}
